%% ma: L11-B16-0041
%% lop: 11
%% chuong: 5
%% bai: 16
%% dang: 11.16.4
%% loai: TN4
%% muc: VD
%% dap_an: "A"
%% nguon: "tai-lieu-hoc-tap-mon-toan-lop-11-theo-chuong-trinh-giao-duc-pho-thong-2018.pdf (Lớp Toán Thầy Duy, Phạm Lê Duy, THPT Chu Văn An), Chương 3 Bài 2, Phần 1, Câu 29 (tr.369)"
%% kien_thuc: "Lớp 11 Bài 16 (giới hạn của hàm số)"
%% trang_thai: chinh-thuc
%% ghi_chu: "Đã duyệt 03/11/2026."
\begin{ex}
Biết $\lim\limits_{x\to+\infty}\dfrac{(2-a)x-3}{x-\sqrt{x^2+1}}=+\infty$ (với $a$ là tham số). Giá trị nhỏ nhất của $P=a^2-2a+4$ là
\choice{\True $4$}{$3$}{$5$}{$1$}
\loigiai{\sloppy $x-\sqrt{x^2+1}=\dfrac{-1}{x+\sqrt{x^2+1}}\to0^-$. Để giới hạn bằng $+\infty$ thì tử phải âm khi $x$ lớn: $2-a<0$, hoặc $2-a=0$ (tử $=-3<0$). Vậy $a\ge2$. Khi đó $P=(a-1)^2+3$ đồng biến trên $[2;+\infty)$, nên $P_{\min}=P(2)=4$.}
\end{ex}
